% !TEX root = main.tex
% Chapter 1: what to do first, the files you edit, how to compile, and what
% to do when compiling fails.
%
% Each chapter lives in its own file, included from main.tex with \input.
% Start the file with \chapter{Title}; type the title in normal
% capitalization (the class sets it in capitals). Headings below it:
%   \section       level 2, Title Case
%   \subsection    level 3, sentence case
%   \subsubsection level 4, sentence case, run-in (the class adds the period)
%   \paragraph     level 5, sentence case, italic (the class adds the period)
% Every heading needs a title; an empty one stops the build with an error.

\chapter{Getting Started}
\label{ch:start}

This chapter tells you what to do first, explains which files you edit, shows how to compile the project locally and on Overleaf, and explains what to do when compiling fails. \cref{ch:requirements} lists the parts an NJIT dissertation must contain, in the order required by the NJIT Graduate Studies Office (GSO), and summarizes the formatting rules the template applies for you. \cref{ch:writing} covers writing the text, \cref{ch:figtables} covers figures and tables, and \cref{ch:references} covers references and the checks to make before you submit. Four appendices follow. \cref{app:appendixA} explains what belongs in an appendix and how to set one up, \cref{app:citestyles} collects guides to the accepted citation styles, \cref{app:packages} describes the packages the template loads, and \cref{app:filetree} lists the template's files.

This template is maintained by Prof.~Roy Goodman and keeps evolving as students use it. If you find a mistake, learn that the GSO has changed a rule, or work out how to handle something the template does not yet cover, please help the students who come after you. Email Prof.~Goodman at \href{mailto:goodman@njit.edu}{\texttt{goodman@njit.edu}}, or make a pull request on GitHub at \url{https://github.com/manroygood/NJIT-Thesis-Template}. Contributors are credited in the template's history.

\section{Start Here}
\label{sec:starthere}

The template comes in two parts. Its main folder holds a short starter document, with two chapters, one appendix, a table, a figure, and a citation, ready for your own text. This sample dissertation, which is also the guide to the template, is in the folder \texttt{sample/}; it uses the class file and settings of the main folder, so it looks exactly like your dissertation will. The sample includes many of the features available in \LaTeX, while the starter contains only a minimal set. Write your dissertation in the main folder, and copy any feature you need from the sample's files into yours. You never need to delete anything from the sample. If you are new to \LaTeX\ or to writing a dissertation, do these five things first. The rest of this document explains each step in more detail.
\begin{enumerate}
\item Compile the starter document, \texttt{main.tex} in the main folder, without changing anything, on Overleaf, in TeXShop, or with \verb+latexmk -pdf main.tex+ (\cref{sec:compiling}). If you get a short PDF with two chapters, a list of references, and an appendix, your installation works.
\item In \texttt{biography.tex}, enter your title, name, degrees, and graduation date. In \texttt{approval.tex}, enter your advisor and committee.
\item Replace the placeholder text of the abstract and the acknowledgments, in the files named after them.
\item Replace the text of \texttt{chapter1.tex}, \texttt{chapter2.tex}, and \texttt{appendix.tex} with your own, and add more chapters as you need them (\cref{sec:setup}).
\item Compile again after every few changes. When an error appears, you then know it comes from the lines you just typed (\cref{sec:errors}).
\end{enumerate}
Words that may be new to you, such as \emph{float}, \emph{package}, and \emph{preamble}, are explained in the Glossary of Terms before this chapter.

\section{Setting up Your Files}
\label{sec:setup}

You edit only a few of the template's files, all in its main folder; \cref{app:filetree} describes every one. To start your own dissertation:
\begin{enumerate}
\item Enter your title, name, degree, department, and graduation date in \texttt{biography.tex}, and your advisor and committee in \texttt{approval.tex}. The class prints them on every page that needs them.
\item Replace the placeholder text of the abstract and the acknowledgments, in the files named after them. For a dedication, copy \texttt{sample/dedication.tex} to the main folder and uncomment its line in \texttt{main.tex}.
\item For a List of Symbols or a Glossary of Terms, copy \texttt{sample/listofsymbols.tex} or \texttt{sample/listofdefinitions.tex} to the main folder, replace the sample entries, and uncomment the matching lines in \texttt{main.tex}.
\item The starter has a single appendix, \texttt{appendix.tex}, printed as APPENDIX with no letter. For two or more appendices, follow \cref{app:appendixA} and this sample's \texttt{main.tex}. Delete the \verb+% !TEX root = main.tex
% SINGLE APPENDIX. Use this file when the document has exactly ONE appendix
% (it is printed as "APPENDIX", with no letter). Start with
% \oneappendix{Title}. For two or more appendices, use appendixA.tex,
% appendixB.tex, appendixC.tex, ... instead (see main.tex). The sample text
% is shared with appendixA.tex through appendix-rules.tex and appendix-demo.tex.
% Every appendix needs a title and a short introduction before any figures
% or tables (GSO). Numbering restarts at A.1.

\oneappendix{Including an appendix or appendices}
\label{app:appendixA}% same label as appendixA.tex, so references in the chapters work either way

% !TEX root = main.tex
% Shared text: the GSO rules for appendices. Included with \input from both
% appendix.tex (single appendix) and appendixA.tex (several appendices), so
% the wording lives in one place. Your own appendices do not need
% this file.

An appendix holds material that supports the dissertation but that a reader does not need in order to follow the argument. A useful test: if a reader skipped the appendix, would the chapters still make sense and still support their conclusions? If not, the material belongs in a chapter.

\textbf{Keep in the chapters} everything a reader needs to follow and evaluate the work:
\begin{itemize}
\item the research questions, the methods, and the main results;
\item the figures and tables that present those results;
\item the key steps of the central proofs and derivations.
\end{itemize}

\textbf{Move to an appendix} material that is long, routine, or of interest mainly to a specialist checking the details:
\begin{itemize}
\item long or routine proofs and derivations, once the chapter has stated the result;
\item full data tables, when the chapter shows a summary;
\item additional figures that repeat a result for other parameters or cases;
\item survey instruments, interview protocols, and consent forms;
\item long source code or pseudocode;
\item details of experimental setups, calibrations, and numerical parameters.
\end{itemize}

Refer to every appendix from the chapter where its material is relevant, for example ``The full proof is given in Appendix~B.'' Material that no chapter refers to probably does not belong in the dissertation at all. Large data sets and complete software packages usually do not belong in an appendix either; upload them to ProQuest as Supplementary Files when you submit the dissertation, or deposit them in a public repository and cite them.

NJIT enforces a strict rule regarding the titling and numbering of appendices:
\begin{itemize}
\item If your dissertation has only a \textbf{single appendix}, it must be titled \textbf{APPENDIX} (without any letter label such as "A"), followed by the appendix title.
\item If your dissertation contains \textbf{two or more appendices}, they must be numbered sequentially as \textbf{APPENDIX A}, \textbf{APPENDIX B}, etc.
\end{itemize}

You should not try to format this manually. Instead:
\begin{itemize}
\item If you have \textbf{one} appendix, use \texttt{appendix.tex} and begin it with the command \verb+\oneappendix{Your Title}+. The heading prints as APPENDIX followed by your title in capitals.
\item If you have \textbf{two or more}, switch to the multi-appendix lines in \texttt{main.tex} and begin each file with \verb+\chapter{Your Title}+. They are lettered A, B, \ldots\ automatically.
\end{itemize}

Every appendix needs a title and a brief introductory description before any equations, figures, or tables (GSO). Appendix pages are numbered continuously with the chapters, and equations, figures, and tables in an appendix are numbered with the appendix letter, starting again at 1 (A.1, B.1, \ldots).


This template uses a single appendix, so its equations, figures, and tables are numbered A.1, A.2, and so on, as below.

% !TEX root = main.tex
% Shared sample: demonstrates numbered equations, figures, and a table inside
% an appendix. Included with \input from appendix.tex (single appendix) and
% appendixA.tex (lettered appendices); only one of those is compiled at a
% time, so each label is defined once. Your own appendices do not
% need this file.

\section{Equations and Figures in Appendices}

Equations inside the appendix are numbered with an appendix prefix:
\begin{align}
\label{eq:app1}
\nabla^2 \phi - \frac{1}{c^2}\frac{\partial^2 \phi}{\partial t^2} &= 0, \\
\label{eq:app2}
\oint_{\partial \Omega} \mathbf{F} \cdot d\mathbf{S} &= \iiint_{\Omega} (\nabla \cdot \mathbf{F})\, dV.
\end{align}
Cross-references like \cref{eq:app1} and \cref{eq:app2} resolve cleanly.

Figures placed within the appendix also receive appropriate appendix numbering, as shown in \cref{fig:app_numeral1} and \cref{fig:app_numeral2}.

\begin{figure}
\centering
\includegraphics[width=2.5in]{drop_cap_1.png}
\caption{An ornamental numeral 1, the first appendix demonstration figure.}
\label{fig:app_numeral1}
\end{figure}

\begin{figure}
\centering
\includegraphics[width=2.5in]{drop_cap_2.png}
\caption{An ornamental numeral 2, the second appendix demonstration figure.}
\label{fig:app_numeral2}
\end{figure}

\section{Tables in an Appendix}

Here is an example of a table in an appendix showing how you might list websites you consulted. \cref{table:latexresources} lists free resources for learning \LaTeX\ and for answering questions as you write. The addresses are printed without the \texttt{https://} prefix, which keeps them short enough to fit the column, but each one still links to the full address:
\begin{verbatim}
\href{https://www.learnlatex.org}{\nolinkurl{www.learnlatex.org}}
\end{verbatim}
For an address printed in full, \verb+\url+ breaks long addresses at any character, thanks to the \texttt{xurl} package.

\begin{table}
\centering
\caption{Free Resources for Learning \LaTeX}
\label{table:latexresources}
\begin{tabularx}{\linewidth}{|>{\raggedright\arraybackslash}p{0.24\linewidth}|>{\raggedright\arraybackslash}p{0.36\linewidth}|>{\raggedright\arraybackslash}X|}
\hline
\textbf{Resource} & \textbf{Address} & \textbf{Good For} \\
\hline
Learn \LaTeX\ in 30 Minutes (Overleaf) & \href{https://www.overleaf.com/learn/latex/Learn_LaTeX_in_30_minutes}{\nolinkurl{www.overleaf.com/learn/latex/Learn_LaTeX_in_30_minutes}} & A first introduction, with examples you can edit online \\
\hline
learnlatex.org & \href{https://www.learnlatex.org}{\nolinkurl{www.learnlatex.org}} & Short interactive lessons covering the basics \\
\hline
Overleaf documentation & \href{https://www.overleaf.com/learn}{\nolinkurl{www.overleaf.com/learn}} & Task-based help on tables, figures, mathematics, and bibliographies \\
\hline
The Not So Short Introduction to \LaTeX & \href{https://ctan.org/pkg/lshort}{\nolinkurl{ctan.org/pkg/lshort}} & A complete free textbook \\
\hline
\LaTeX\ Wikibook & \href{https://en.wikibooks.org/wiki/LaTeX}{\nolinkurl{en.wikibooks.org/wiki/LaTeX}} & A searchable reference organized by topic \\
\hline
How Should We \LaTeX? (Roy Goodman) & \href{https://roygoodman.net/post/how-to-latex/}{\nolinkurl{roygoodman.net/post/how-to-latex}} & Style habits that manuals skip: meaningful macros, consistent notation, punctuating equations, and coauthoring \\
\hline
Short Math Guide for \LaTeX\ (AMS) & \href{https://ctan.org/pkg/short-math-guide}{\nolinkurl{ctan.org/pkg/short-math-guide}} & The commands of \texttt{amsmath} at a glance \\
\hline
Detexify & \href{https://detexify.kirelabs.org}{\nolinkurl{detexify.kirelabs.org}} & Finding the command for a symbol by drawing it \\
\hline
Comprehensive \LaTeX\ Symbol List & \href{https://ctan.org/pkg/comprehensive}{\nolinkurl{ctan.org/pkg/comprehensive}} & Every symbol available, with the package that provides it \\
\hline
TeX Stack Exchange & \href{https://tex.stackexchange.com}{\nolinkurl{tex.stackexchange.com}} & Answers to specific questions; search before asking \\
\hline
CTAN & \href{https://ctan.org}{\nolinkurl{ctan.org}} & Documentation for every package \\
\hline
\end{tabularx}
\end{table}

% \clearpage prints the floating Table A.1 first, so the tables stay in
% order; a longtable does not float and would otherwise come out first.
\clearpage
\section{Tables Longer Than One Page}

A table too long for one page is typed with the \texttt{longtable} environment instead of a \texttt{table} environment, as in \cref{table:longsample}. It does not float: it starts where it is typed and continues onto as many pages as it needs. Its title goes in the \verb+\endfirsthead+ part; the \verb+\endhead+ part, which starts with the class command \verb+\tablecontinued+, is repeated at the top of every later page, so each of them begins with ``\textbf{Table X.x (Continued)}'' and the column headings (GSO). The continued title is not listed again in the List of Tables.

\begin{longtable}{>{\centering\arraybackslash}p{0.15\linewidth}l>{\raggedright\arraybackslash}p{0.45\linewidth}}
\caption{Commands for Common Mathematical Symbols}
\label{table:longsample} \\
\toprule
\textbf{Symbol} & \textbf{Command} & \textbf{Meaning} \\
\midrule
\endfirsthead
\tablecontinued \\
\toprule
\textbf{Symbol} & \textbf{Command} & \textbf{Meaning} \\
\midrule
\endhead
\bottomrule
\endlastfoot
$\alpha$ & \verb+\alpha+ & Lowercase alpha \\
$\beta$ & \verb+\beta+ & Lowercase beta \\
$\gamma$ & \verb+\gamma+ & Lowercase gamma \\
$\delta$ & \verb+\delta+ & Lowercase delta \\
$\epsilon$ & \verb+\epsilon+ & Lowercase epsilon \\
$\zeta$ & \verb+\zeta+ & Lowercase zeta \\
$\eta$ & \verb+\eta+ & Lowercase eta \\
$\theta$ & \verb+\theta+ & Lowercase theta \\
$\iota$ & \verb+\iota+ & Lowercase iota \\
$\kappa$ & \verb+\kappa+ & Lowercase kappa \\
$\lambda$ & \verb+\lambda+ & Lowercase lambda \\
$\mu$ & \verb+\mu+ & Lowercase mu \\
$\nu$ & \verb+\nu+ & Lowercase nu \\
$\xi$ & \verb+\xi+ & Lowercase xi \\
$\pi$ & \verb+\pi+ & Lowercase pi \\
$\rho$ & \verb+\rho+ & Lowercase rho \\
$\sigma$ & \verb+\sigma+ & Lowercase sigma \\
$\tau$ & \verb+\tau+ & Lowercase tau \\
$\upsilon$ & \verb+\upsilon+ & Lowercase upsilon \\
$\phi$ & \verb+\phi+ & Lowercase phi \\
$\chi$ & \verb+\chi+ & Lowercase chi \\
$\psi$ & \verb+\psi+ & Lowercase psi \\
$\omega$ & \verb+\omega+ & Lowercase omega \\
$\varepsilon$ & \verb+\varepsilon+ & Variant of epsilon \\
$\vartheta$ & \verb+\vartheta+ & Variant of theta \\
$\varpi$ & \verb+\varpi+ & Variant of pi \\
$\varphi$ & \verb+\varphi+ & Variant of phi \\
$\Gamma$ & \verb+\Gamma+ & Uppercase gamma \\
$\Delta$ & \verb+\Delta+ & Uppercase delta \\
$\Theta$ & \verb+\Theta+ & Uppercase theta \\
$\Lambda$ & \verb+\Lambda+ & Uppercase lambda \\
$\Xi$ & \verb+\Xi+ & Uppercase xi \\
$\Pi$ & \verb+\Pi+ & Uppercase pi \\
$\Sigma$ & \verb+\Sigma+ & Uppercase sigma \\
$\Upsilon$ & \verb+\Upsilon+ & Uppercase upsilon \\
$\Phi$ & \verb+\Phi+ & Uppercase phi \\
$\Psi$ & \verb+\Psi+ & Uppercase psi \\
$\Omega$ & \verb+\Omega+ & Uppercase omega \\
$\nabla$ & \verb+\nabla+ & Gradient (nabla) \\
$\partial$ & \verb+\partial+ & Partial derivative \\
$\infty$ & \verb+\infty+ & Infinity \\
$\sum$ & \verb+\sum+ & Summation \\
$\prod$ & \verb+\prod+ & Product \\
$\int$ & \verb+\int+ & Integral \\
$\leq$ & \verb+\leq+ & Less than or equal to \\
$\geq$ & \verb+\geq+ & Greater than or equal to \\
$\neq$ & \verb+\neq+ & Not equal to \\
$\approx$ & \verb+\approx+ & Approximately equal to \\
$\in$ & \verb+\in+ & Element of \\
$\forall$ & \verb+\forall+ & For all \\
$\to$ & \verb+\to+ & Maps to, tends to \\
$\times$ & \verb+\times+ & Times, cross product \\
\end{longtable}

+ line if you have no appendix.
\item Write each chapter in its own file, \texttt{chapter1.tex}, \texttt{chapter2.tex}, and so on, and list them with \verb+\input+ in \texttt{main.tex}. While drafting, you can comment out finished chapters to compile faster. Start every new \texttt{.tex} file with the line \verb+% !TEX root = main.tex+ (\cref{subsec:editors}).
\item Put your sources in \texttt{references.bib} and your images in the \texttt{Graphics/} folder of the main folder. Replace the starter's placeholder entries in \texttt{references.bib}; this sample's \texttt{sample/references.bib} has a model of each common entry type.
\item Set the citation style for your discipline with the \texttt{style=} option on the \texttt{biblatex} line in \texttt{packages.tex}. \cref{tab:phd-styles} in \cref{sec:citestyles} lists the styles each program recommends, and \cref{ch:references} explains how to cite sources and build the bibliography.
\item Load any packages you need in \texttt{packages.tex}, before the \texttt{hyperref} line, and define your own notation and theorem environments in \texttt{macros.tex}.
\end{enumerate}
For a master's thesis, add the \texttt{thesis} option to the \verb+\documentclass+ line at the top of \texttt{main.tex}:
\begin{verbatim}
\documentclass[12pt,thesis]{njit_thesis}
\end{verbatim}
The title page then reads ``A Thesis'', the advisor is listed as ``Thesis Advisor'' on the approval page, and the copyright page is replaced by the blank page the GSO requires. In \texttt{biography.tex}, change \verb+\thisdegree+ to the full name of your master's degree, and if your degree is not a Master of Science, set the ``Degree:'' line of the Biographical Sketch with \verb+\degreelevel+. Two limits also change: a thesis abstract may be only one page long, where a dissertation abstract may be two, and a thesis committee needs at least three members, including the advisor, where a dissertation committee needs five. You should not need to edit the class file, \texttt{njit\_thesis.cls}.

To compile this sample, run \verb+latexmk -pdf main.tex+ in the \texttt{sample/} folder, or open \texttt{sample/main.tex} in TeXShop. It finds the class file, \texttt{packages.tex}, \texttt{macros.tex}, and the citation styles in the main folder through the paths \verb+../njit_thesis+ and \verb+../packages+. On Overleaf, read the sample as a PDF from the template's GitHub Releases page instead, because Overleaf expects the main file of a project in its top folder.

Keep your work safe from the first day. A lost laptop or an overwritten file a month before the deadline is a common disaster, and an easy one to prevent. Keep the folder under version control with Git (the free GitHub Desktop application is enough), or rely on Overleaf's history, and keep a second copy somewhere other than your own computer.

\section{Compiling the Dissertation}
\label{sec:compiling}
The subsections below describe several ways to compile. If compiling stops with an error that a \texttt{.sty} file was not found, a package the template needs is not installed; \cref{sec:missingpackages} explains how to find and install it.


\subsection{Recommended: \texttt{latexmk}}
The easiest and most reliable way to compile your dissertation locally is with \texttt{latexmk}:
\begin{verbatim}
latexmk -pdf main.tex
\end{verbatim}
\texttt{latexmk} monitors all file dependencies, automatically invokes \texttt{biber} for the bibliography and publications list, and runs \texttt{pdflatex} as many times as necessary to ensure all cross-references, Table of Contents entries, and list page numbers are completely up to date.

\subsection{Manual compilation steps}
If running commands manually in a terminal:
\begin{enumerate}
  \item \verb+pdflatex main.tex+ (generates auxiliary files and bibliography requests)
  \item \verb+biber main+ (processes \texttt{references.bib} and resolves citations)
  \item \verb+pdflatex main.tex+ (incorporates bibliography and generates cross-references)
  \item \verb+pdflatex main.tex+ (resolves final page numbers for TOC and lists)
\end{enumerate}

\subsection{Using Overleaf}
The template works on Overleaf without changes:
\begin{enumerate}
\item Download the template as a ZIP file and upload it to Overleaf with \emph{New Project}, then \emph{Upload Project}. You may delete the \texttt{sample/} folder from the project first; Overleaf compiles only the starter document.
\item In the project's \emph{Menu}, set the compiler to \textbf{pdfLaTeX} and the main document to \texttt{main.tex}.
\item Overleaf normally detects that the bibliography uses Biber. If the references show as ``??'', look for a bibliography setting in \emph{Menu} and choose \textbf{Biber}.
\end{enumerate}

\subsection{Using TeXShop, TeXstudio, or VS Code}
\label{subsec:editors}
Every \texttt{.tex} file in the template begins with the line
\begin{verbatim}
% !TEX root = main.tex
\end{verbatim}
This comment tells TeXShop on a Mac, and several other editors, to compile \texttt{main.tex} whichever file is open, so you can compile while editing a chapter. Keep it in every file, including new ones you create. In \texttt{main.tex}, the line \verb+% !TEX TS-program = pdflatexmk+ tells TeXShop to compile with \texttt{latexmk}, which runs Biber when needed. In other editors, choose pdf\LaTeX\ as the engine and Biber as the bibliography tool, or choose \texttt{latexmk}: in TeXstudio, under \emph{Options}, \emph{Configure TeXstudio}, \emph{Build}; in VS Code, with the LaTeX Workshop extension, whose default recipe uses \texttt{latexmk}.

\section{When Compiling Fails}
\label{sec:errors}
When \LaTeX\ meets something it cannot handle, it stops and prints an error that begins with an exclamation mark, followed by the line where it noticed the problem:
\begin{verbatim}
! Undefined control sequence.
l.42 The value of \alpah
\end{verbatim}
Read the first error only; later errors are often caused by the first one. The number after \texttt{l.} is the line in the file being read, and the error is usually on that line or just before it. If the compiler stops with a \texttt{?} prompt, type \texttt{X} and press Enter to quit, fix the line, and compile again. The most common errors are:
\begin{itemize}
\item \textbf{Undefined control sequence}: A command is misspelled, as in \verb+\alpah+ above, or comes from a package that is not loaded.
\item \textbf{Missing \$ inserted}: A math command such as \verb+^+, \verb+_+, or \verb+\alpha+ is used outside math mode, or a \verb+$+ is missing.
\item \textbf{Missing \} inserted, or Runaway argument}: A brace \verb+{+ is not closed. Check the lines just before the one reported.
\item \textbf{File not found}: A package or image is missing, or its name is misspelled. For packages, see \cref{sec:missingpackages}; for images, check the name and the \texttt{Graphics/} folder.
\item \textbf{Citation or reference undefined (a warning, shown as ``??'')}: Compile again; if it persists, the key in \verb+\cite+ or \verb+\ref+ does not match any \verb+\label+ or \texttt{.bib} entry.
\end{itemize}
If the PDF does not change or the errors make no sense, delete the auxiliary files (\cref{app:filetree} lists them) and compile again. If you are stuck, search for the exact error message; TeX Stack Exchange (\url{https://tex.stackexchange.com}) has answers to almost all of them.
